\documentclass[11pt]{article}
\usepackage[margin=1in]{geometry}
\usepackage{amsmath,amssymb,bm,booktabs,array,longtable}
\usepackage[hidelinks]{hyperref}
\usepackage{xcolor,listings}
\lstset{basicstyle=\ttfamily\small,breaklines=true,frame=single,
  keywordstyle=\color{blue!60!black},commentstyle=\color{green!40!black}}
\title{Phase-by-Phase FNO and PINO Training for WaveQLab2D}
\author{WaveQLab2D Simulation Pipeline}
\date{\today}

\begin{document}
\maketitle
\tableofcontents

\section{Purpose and experimental contract}
The objective is to learn the map from a Gaussian source position and query time to
the full particle-velocity field,
\[
  \mathcal{G}: (x_0,y_0,t)\longmapsto
  \bm v(x,y,t)=(v_x(x,y,t),v_y(x,y,t)),
\]
on the fixed $10\times10$ km objective window.  The supervised FNO and PINO use the
same raw simulations, prepared tensors, normalization, architecture, optimizer
family, validation cases, and test cases.  PINO is warm-started from FNO and adds a
wave-equation residual.  This isolates the effect of physics regularization.

The pilot contains 500 simulations: 350 training, 50 validation, 50 ID test, and 50
far-OOD test cases.  Training source offsets are in $[-3,3]$ km; far-OOD offsets are
in $[-5,-4)\cup(4,5]$ km.  Designs at 500, 1000, 1500, and 2000 samples are nested,
so learning curves do not change earlier examples.

\section{Governing equations}
For homogeneous isotropic linear elasticity,
\begin{align}
 \rho\,\partial_t \bm v &= \nabla\cdot\bm\sigma + \bm f,\\
 \partial_t\bm\sigma &= \lambda(\nabla\cdot\bm v)\bm I
        +\mu\left(\nabla\bm v+\nabla\bm v^T\right),
\end{align}
where $\rho$ is density and $\lambda,\mu$ are Lam\'e parameters.  Eliminating stress
away from the source gives the velocity Navier equation
\begin{equation}
 \partial_{tt}\bm v = c_s^2\nabla^2\bm v
 +(c_p^2-c_s^2)\nabla(\nabla\cdot\bm v),
 \label{eq:navier}
\end{equation}
with $c_p^2=(\lambda+2\mu)/\rho$ and $c_s^2=\mu/\rho$.

In components,
\begin{align}
R_x &= v_{x,tt}-c_p^2v_{x,xx}-c_s^2v_{x,yy}
 -(c_p^2-c_s^2)v_{y,xy},\\
R_y &= v_{y,tt}-c_s^2v_{y,xx}-c_p^2v_{y,yy}
 -(c_p^2-c_s^2)v_{x,xy}.
\end{align}
The PINO physics loss is the mean squared nondimensional residual outside a source
mask.  Centered second differences use the saved spacing $\Delta t_s=0.02$ s and
$\Delta x=\Delta y=0.1$ km.  The source mask radius is 0.5 km, exceeding twice the
solver Gaussian standard deviation $\sigma=2\Delta x=0.2$ km.

\section{Fourier Neural Operator}
Let $a(x,y)$ denote the input-channel field.  A spectral layer applies
\begin{equation}
 z_{\ell+1}=\phi\!\left(W_\ell z_\ell+
 \mathcal{F}^{-1}\!\left(R_\ell\,\mathcal{F}z_\ell\right)\right),
\end{equation}
where $W_\ell$ is a learned pointwise map and $R_\ell$ contains learned complex
weights for a truncated set of Fourier modes.  The retained low modes efficiently
represent propagating wave structure while the local $1\times1$ convolution restores
local coupling.

The model inputs are: the spatial Gaussian source, source multiplied by its temporal
envelope, normalized $x$ and $y$, distance to the objective boundary, normalized
query time, normalized $y_0$, and three sine/cosine Fourier pairs in time.  There are
13 channels and two output channels.

\section{Losses}
The supervised loss combines pointwise, gradient, and spectral agreement:
\begin{equation}
 \mathcal{L}_{\rm data}=\operatorname{MSE}(\hat v,v)
 +0.05\operatorname{MAE}(\hat v,v)
 +0.05(\operatorname{MAE}(D_x\hat v,D_xv)+\operatorname{MAE}(D_y\hat v,D_yv))
 +0.01\operatorname{MAE}(|\mathcal{F}\hat v|,|\mathcal{F}v|).
\end{equation}
PINO minimizes
\begin{equation}
 \mathcal{L}_{\rm PINO}=\mathcal{L}_{\rm data}+\gamma\mathcal{L}_{\rm physics},
 \qquad
 \mathcal{L}_{\rm physics}=\frac{1}{|\Omega_m|}\sum_{\Omega_m}
 \left\|\frac{\bm R}{v_{\rm scale}/\Delta t_s^2}\right\|_2^2.
\end{equation}
The default $\gamma=0.01$ follows the earlier pilot.  It must be tuned only on the
validation split; recommended values are $10^{-3},3\times10^{-3},10^{-2},3\times
10^{-2}$.

\section{Phase algorithms}
\subsection{Phase 01: integrity and preparation}
\texttt{phase\_01.py} finds exactly one finalized raw HDF5 archive per manifest row,
merges the left/right blocks while removing the duplicate interface point, selects
only $v_x,v_y$, transposes to $[t,x,y,2]$, checks shape $[250,101,101,2]$ and finite
values, then writes a chunked compressed prepared HDF5 file.
\begin{lstlisting}[language=Python]
left = source["DomainOutput_l"][:-1, :, :, :2]
right = source["DomainOutput_r"][:, :, :, :2]
velocity = np.concatenate((left, right), axis=0).transpose(2,0,1,3)
\end{lstlisting}

\subsection{Phase 02: leakage-free normalization}
\texttt{phase\_02.py} reads training cases only.  Per-component maximum magnitude and
RMS are accumulated in time chunks.  The scale is
\[s_k=\max(\max|v_k|,5\operatorname{RMS}(v_k)).\]
Coordinates, saved times, spacings, scales, and field order are stored in
\texttt{normalization.json}.  Validation and test data never influence this file.

\subsection{Phase 03: baseline FNO}
Each optimization step samples cases and query frames, constructs the conditioning
channels, scales the target velocities, evaluates the FNO, and applies the composite
supervised loss.  AdamW, cosine learning-rate decay, unit gradient clipping, and
validation relative $L_2$ checkpoint selection are used.

\subsection{Phase 04: bounded hyperparameter tuning}
\texttt{phase\_04.py} searches width $\{24,32,48\}$, modes $\{12,16,20\}$, and
learning rate $\{5\times10^{-4},10^{-3}\}$.  Short trials rank configurations by
validation relative $L_2$.  The winning configuration should then be retrained with
the full Phase 03 budget and at least three seeds.  The test set is not consulted.

\subsection{Phase 05: PINO fine-tuning}
Three consecutive query times are passed through the same FNO.  Their predicted
fields approximate $v_{tt}$, while spatial differences approximate the derivatives
in Eq.~\eqref{eq:navier}.  Fine-tuning starts from the best FNO checkpoint with a
smaller learning rate $10^{-4}$ to preserve the learned data manifold.
\begin{lstlisting}[language=Python]
predicted = model(triplet_inputs)
data = supervised_loss(predicted_current, truth_current)
physics = navier(predicted, cases, metadata, scale)
loss = data + physics_weight * physics
\end{lstlisting}

\subsection{Phase 06: matched evaluation}
Both checkpoints predict every saved test frame.  The report separates ID and
far-OOD cases and records global relative $L_2$, per-case relative $L_2$, maximum
absolute error, median/p90 error, total inference time, and milliseconds per frame.
Evaluation uses identical batches and normalization for FNO and PINO.

\section{Parameter and hyperparameter rationale}
\begin{longtable}{p{0.25\linewidth}p{0.2\linewidth}p{0.46\linewidth}}
\toprule Parameter & Default & Rationale\\\midrule
Spatial modes & 16 & Resolves wavelengths represented on a $101^2$ grid without the cost of retaining all modes.\\
Width & 32 & Moderate H100 memory use; width search tests under/over-capacity.\\
Layers & 4 & Repeated global/local mixing with stable optimization.\\
Batch size & 4 FNO, 2 PINO & PINO evaluates three times per sampled case.\\
FNO learning rate & $10^{-3}$ & AdamW baseline; searched against $5\times10^{-4}$.\\
PINO learning rate & $10^{-4}$ & Conservative warm-start fine-tuning.\\
Weight decay & $10^{-4}$ & Limits spectral-weight overfit.\\
Physics weight & $10^{-2}$ & Previous pilot baseline; tune logarithmically.\\
Source mask & 0.5 km & Excludes the forced region where homogeneous Eq.~\eqref{eq:navier} is invalid.\\
Checkpoint metric & validation relative $L_2$ & Directly targets physical field accuracy.\\
\bottomrule
\end{longtable}

\section{Accuracy-first tuning protocol}
First verify overfitting on a tiny subset; failure indicates a code or conditioning
problem.  Next tune FNO capacity and optimizer settings using all validation cases.
Retrain the winner with three seeds.  Tune PINO physics weight from weak to strong,
rejecting settings that reduce residual but degrade validation field error.  Finally
evaluate once on ID and far-OOD test groups.  For the nested 500/1000/1500/2000
study, preserve hyperparameter-search rules and report both single-seed and ensemble
learning curves.

\section{Reproducible execution}
\begin{lstlisting}[language=bash]
python phase_01.py DATASET
python phase_02.py DATASET
python phase_03.py DATASET --output-dir DATASET/models/fno
python phase_04.py DATASET --output-dir DATASET/models/search
python phase_05.py DATASET DATASET/models/fno/best.pt \
  --output-dir DATASET/models/pino
python phase_06.py DATASET --fno DATASET/models/fno/best.pt \
  --pino DATASET/models/pino/best.pt --output-dir DATASET/evaluation
\end{lstlisting}

\end{document}
